\chapter{Bleeding Gums}

\section*{Definition}
Bleeding gums means bleeding from the gum margin, often noticed during brushing, flossing, or eating. It is commonly caused by plaque-related gingivitis, but persistent, spontaneous, or disproportionate bleeding can have other dental or medical causes.

\section*{When to See a Doctor}

\begin{itemize}
 \item Bleeding that is spontaneous, profuse, or difficult to control with direct pressure
 \item Gingival bleeding accompanied by petechiae, ecchymoses, nosebleed, or prolonged bleeding from minor cuts (suggests a systemic bleeding disorder)
 \item Bleeding with gingival hyperplasia, ulceration, or necrotic papillae (ANUG, leukemia, scurvy)
 \item Fever, malaise, swollen neck glands, unexplained bruising, or weight loss
 \item Any bleeding when brushing or eating should be discussed with a dentist, particularly if it is recurrent or persists despite improved oral hygiene.
 \item Unexplained gingival bleeding in the setting of liver disease or new anticoagulation therapy
\end{itemize}

\section*{How It Develops}
Plaque that is not removed irritates the gum margin and triggers inflammation. Inflamed gums become red, swollen, and more likely to bleed when touched. If inflammation progresses to periodontitis, the tissues and bone supporting the teeth can be damaged; gingivitis itself is generally reversible with effective plaque control.

\section*{Causes}
\begin{itemize}
 \item Local causes: Gingivitis or periodontitis, plaque and calculus, food impaction, irritation from dental appliances, and acute necrotizing gingivitis
 \item Systemic bleeding disorders: Immune thrombocytopenia (ITP), von Willebrand disease, hemophilia A/B (Factor VIII/IX deficiency), platelet function disorders, liver disease–related coagulopathy (impaired vitamin K–dependent factor synthesis), vitamin K deficiency
 \item Medication-induced: Anticoagulants (warfarin, heparin, direct oral anticoagulants such as apixaban and rivaroxaban), antiplatelet agents (aspirin, clopidogrel), certain antihypertensives (calcium channel blockers like nifedipine, amlodipine can cause gingival hyperplasia with secondary bleeding), methotrexate, valproic acid (thrombocytopenia), chemotherapy
 \item Nutritional deficiencies: Vitamin C deficiency (scurvy—gingival hypertrophy, hemorrhage, and loose teeth with impaired collagen synthesis), vitamin K deficiency
 \item Hematologic malignancy: Acute leukemia (especially acute monocytic or acute myeloid leukemia with gingival infiltration), multiple myeloma, myelodysplastic syndromes
 \item Hormonal: Pregnancy gingivitis (estrogen and progesterone increase gingival vascularity), puberty gingivitis, oral contraceptive use
 \item Systemic inflammation: Granulomatosis with polyangiitis (Wegeners), sarcoidosis, Crohn disease (oral mucosal involvement), systemic lupus erythematosus
 \item Infectious: Acute herpetic gingivostomatitis (HSV-1), acute necrotizing ulcerative gingivitis (fusospirochetal infection), candidiasis, HIV-associated gingivitis and linear gingival erythema
 \item Neoplastic: Oral squamous cell carcinoma presenting as a friable, bleeding lesion; pyogenic granuloma; Kaposi sarcoma (HHV-8)
\end{itemize}

\section*{Related Symptoms}
\begin{itemize}
 \item Gingival erythema
 \item Gingival edema and hypertrophy
 \item Gingival tenderness or pain
 \item Halitosis
 \item Tooth mobility
 \item Gingival recession
 \item Oral mucosal petechiae
 \item nosebleed
 \item Easy bruising
 \item Prolonged bleeding from wounds
\end{itemize}
\textbf{Important:} While bleeding gums most often reflect plaque-induced gingivitis, new-onset or disproportionate bleeding — especially with systemic symptoms — should prompt evaluation for thrombocytopenia, coagulopathy, leukemia, or vitamin C deficiency.

\section*{Medical Evaluation}
\begin{itemize}
 \item History includes bleeding pattern (provoked vs. spontaneous), duration, oral hygiene habits, medication list (especially anticoagulants and antiplatelets), alcohol use, and systemic symptoms.
 \item Full intraoral examination with periodontal probing to assess pocket depth, bleeding on probing (BOP score), attachment loss, and presence of calculus or defective restorations.
 \item Blood tests are not routine for typical gingivitis. For spontaneous or disproportionate bleeding or systemic symptoms, a clinician may order a complete blood count and coagulation tests, with further tests guided by the findings.
 \item Nutritional assessment and vitamin testing are considered when the history suggests a deficiency; supplements should not be started at high doses without advice.
\end{itemize}

\section*{Treatment Options}
\begin{itemize}
 \item Professional dental cleaning (scaling and root planing) to remove plaque and calculus; instruction in proper brushing and flossing technique.
 \item An antimicrobial mouthrinse may be recommended for selected patients as an adjunct, but it does not replace brushing and interdental cleaning. Chlorhexidine can stain teeth and alter taste and should be used only as directed.
 \item Address underlying systemic cause: Vitamin C or K repletion, adjustment of anticoagulation under medical supervision, treatment of leukemia or bleeding diathesis, withdrawal or substitution of causative medications when possible.
 \item Surgical intervention: Gingivectomy or gingivoplasty for drug-induced or genetic gingival hyperplasia after medical management; biopsy of persistent, atypical lesions to exclude malignancy.
 \item Topical haemostatic agents (absorbable gelatin sponge, oxidized cellulose, or tranexamic-acid mouthwash) may be used by clinicians for acute bleeding in someone with coagulopathy.
\end{itemize}

\section*{Self-Care and Home Management}
\begin{itemize}
 \item Brush gently twice daily with fluoride toothpaste and clean between the teeth daily with floss or an interdental brush recommended for you. Do not stop cleaning solely because the gums bleed.
 \item Use a warm saline rinse (1/2 teaspoon salt in 8 oz water) after brushing to soothe inflamed gums.
 \item Avoid vigorous brushing, hard-bristled toothbrushes, and abrasive toothpaste during active bleeding episodes.
 \item Discontinue use of tobacco products and limit alcohol intake, both of which exacerbate gingival inflammation.
 \item If anticoagulated, confirm with the prescribing physician before modifying medication; iced water held in the mouth may reduce active oozing.
\end{itemize}

\section*{References}
\begin{thebibliography}{99}
\bibitem{bleeding_gums:ref1} Tonetti MS, Greenwell H, et al. ``Staging and grading of periodontitis: framework and proposal.'' \textit{J Periodontol}. 2018;89(Suppl 1):S159-S172.
\bibitem{bleeding_gums:ref2} Chapple ILC, Mealey BL, et al. ``Periodontal health and gingival diseases.'' \textit{J Periodontol}. 2018;89(Suppl 1):S74-S84.
\bibitem{bleeding_gums:ref3} Newman MG, Takei HH, et al. \textit{Newman and Carranza\'s Clinical Periodontology}, 13th ed. Elsevier, 2018.
\bibitem{bleeding_gums:ref4} Pihlstrom BL, Michalowicz BS, et al. ``Periodontal diseases.'' \textit{Lancet}. 2005;366(9499):1809-1820.
\bibitem{bleeding_gums:ref5} National Health Service. ``Gum disease.'' Reviewed 2025. https://www.nhs.uk/conditions/gum-disease/.
\bibitem{bleeding_gums:ref6} American Dental Association. ``Floss/Manual Interdental Cleaners.'' Accessed September 2026. https://www.ada.org/resources/ada-library/oral-health-topics/floss.
\end{thebibliography}
